\documentclass{article}
\usepackage{graphicx}
\usepackage{amsmath}
\usepackage{amssymb}
\usepackage{amsthm}
\usepackage{tabularx}
\usepackage{booktabs}
\usepackage{multirow}
\usepackage{stmaryrd}

\newtheorem{definition}{Definition}

\title{PP1 - Vyhodnocování úplnosti testů vůči softwarovým požadavkům - formalismus}
\author{  }
\date{April 2026}

\begin{document}

\maketitle

\section{The primitives}

These primitive sets serve as the base that the rest of the formalism builds on. The primitive sets are the following:
\begin{itemize}
    \item $T_D$, $T_C$ -- the \textit{discrete time} and \textit{continuous time} sets.
          An infinite totally ordered set; the time set can have \textit{discrete} or \textit{continuous} flavour; exactly one flavour is chosen per requirement.
    \item $V$ -- the $Values$ set. These can be used to describe the values of entities.
\end{itemize}

\begin{definition}[Discrete time set]\label{def:time_discrete}
    The \emph{discrete time set} $T_D$ is totally ordered and countable.
    \[
        T_D = \mathbb{N}
    \]
\end{definition}

\begin{definition}[Continuous time set]\label{def:time_continuous}
    The \emph{continuous time set} $T_C$ is totally ordered and uncountable.
    \[
        T_C = \mathbb{R}_{\geq 0}
    \]
\end{definition}

\begin{definition}[Duration]\label{def:duration}
    A \emph{duration} is a non-negative element of the time set, representing the length of a time span:
    \[
        \Delta = T
    \]
    Thus $\Delta = \mathbb{N}$ for discrete time and $\Delta = \mathbb{R}_{\geq 0}$ for continuous time.
\end{definition}

\begin{definition}[Interval]\label{def:interval}
    An \emph{interval} is a pair of a start time point and a duration:
    \[
        \mathcal{I} = T \times \Delta
    \]
    An interval $(t, d) \in \mathcal{I}$ represents the time span $[t,\, t + d]$.
\end{definition}

\begin{definition}[Values]\label{def:values}
    The \emph{values set} $V$ is the union of the following sets:
    \[
        V = \mathbb{B} \cup \mathbb{R} \cup T \cup \mathcal{I}
    \]
    where $\mathbb{B} = \{\mathit{true}, \mathit{false}\}$ is the set of booleans, $\mathbb{R}$ is the set of real numbers, $T$ is the time set, and $\mathcal{I}$ is the set of intervals.
\end{definition}

\section{Entities and entity types}

An entity type is an item belonging to a set of predefined entity types.

\begin{definition}[Entity type set]\label{def:entity_type_set}
    \[
        T_\text{E} = \{ \text{Signal}, \text{Storage}, \text{Event}, \text{Channel}, \text{State}, \text{Value}, \text{Abstract}\}
    \]
\end{definition}

\begin{definition}[Entity name set]\label{def:entity_name_set}
    \[
        N_\text{E} \subseteq \Sigma^*
    \]

    , where $\Sigma$ is an alphabet.
\end{definition}

Entity type modifiers are properties which can be attached to
an entity of a particular type. Semantics are defined for each modifier.

\begin{definition}[Entity type modifiers set]\label{def:entity_type_modifiers_set}
    \[
        M = \{m\,|\,m\,\text{is a modifier of an entity type}\}
    \]

\end{definition}

Each entity type is compatible with only a subset of the entity modifiers.
Therefore we define the \textit{entity type modifiers function}.

\begin{definition}[Entity type modifiers function]\label{def:entity_type_modifiers_fn}
    \[
        M_\text{E}: T_E \rightarrow \mathcal{P}(M)
    \]

    , where $\mathcal{P}(M)$ is the power set of $M$.
\end{definition}

% \begin{definition}[Entity type]\label{def:entity_type}
%     \[
%         T_e
%     \]
% \end{definition}

An entity is something that is referenced in the requirement.
Formally, entity is defined as a four-item tuple.

\begin{definition}[Entity]\label{def:entity}
    \[
        e = (n_e, t_e, m_e, d_e)
    \]

    , where
    \begin{itemize}
        \item $n_e$ is the name of the entity; $n_e \in N_E$
        \item $t_e$ is the name of the entity; $t_e \in T_E$
        \item $m_e$ is the set of modifiers of the entity; $m_e \in T_E$, $m_e \subseteq M_E(t_e)$
        \item $d_e$ is the description of the entity; $d_e \in \Sigma^*$
    \end{itemize}
\end{definition}

\section{Valuations and traces}

\begin{definition}[Valuation]\label{def:valuation}
    A valuation is a function $\sigma : E \to V$ assigning a value to every entity.
    For entities of type \texttt{Event}, $\sigma(e) \in \{\mathit{true}, \mathit{false}\}$,
    where $\sigma(e) = \mathit{true}$ iff the event occurs at that time point.
\end{definition}

We write $\Sigma$ for the set of all valuations.

\begin{definition}[Trace]\label{def:trace}
    A trace is a function $\rho : T \to \Sigma$. We write $\rho(t)(e)$ for the value of entity $e$ at time $t$.
\end{definition}

\section{Expressions and predicates}

\subsection{Expressions}

Expressions are terms that evaluate to a value in $V$ given a trace $\rho$.

\begin{definition}[]\label{def:expr_val}
    For a trace $\rho \in \Sigma^T$, the function $\mathit{Val}_\rho: E \times T \rightarrow V$ denotes the value of entity $e \in E$ at time $t \in T$:
    \[
        \mathit{Val}_\rho(e,\, t) \;=\; \rho(t)(e)
    \]
\end{definition}

\begin{definition}[Event occurrence count]\label{def:evt_occ_count}
    For a trace $\rho \in \Sigma^T$ and an entity $ev \in E$ of type \texttt{Event},
    $\mathit{EvtOccCount}_\rho : E \times \mathcal{I} \rightarrow \mathbb{N}$ returns the number of occurrences of $ev$ within an interval $i \in \mathcal{I}$:
    \[
        \mathit{EvtOccCount}_\rho(ev,\, i) \;=\; \bigl|\{\, s \in T \mid \mathit{Start}(i) \leq s \leq \mathit{End}(i) \;\land\; \mathit{Val}_\rho(ev,\, s) = \mathit{true} \,\}\bigr|
    \]
\end{definition}

\begin{definition}[]\label{def:last_occ}
    For a trace $\rho \in \Sigma^T$ and an entity $ev \in E$ of type \texttt{Event},
    $\mathit{LastOcc}_\rho : E \times T \rightarrow T \cup \{-\infty\}$ returns the time of the most recent occurrence of $ev$ up to time $t$:
    \[
        \mathit{LastOcc}_\rho(ev,\, t) \;=\;
        \begin{cases}
            \max\{\, s \in T \mid s \leq t \;\land\; \mathit{Val}_\rho(ev,\, s) = \mathit{true} \,\} & \text{if the set is non-empty} \\
            -\infty                                                                                  & \text{otherwise}
        \end{cases}
    \]
\end{definition}

\begin{definition}[Maximum value function]\label{def:max_val}
    For a trace $\rho \in \Sigma^T$ and an entity $e \in E$ with an ordered value domain,
    $\mathit{MaxVal}_\rho : E \times T \rightarrow V$ returns the maximum value of $e$ up to and including time $t$:
    \[
        \mathit{MaxVal}_\rho(e,\, t) \;=\; \sup\{\, \mathit{Val}_\rho(e,\, s) \mid s \in T,\; s \leq t \,\}
    \]
\end{definition}

\begin{definition}[Minimum value function]\label{def:min_val}
    For a trace $\rho \in \Sigma^T$ and an entity $e \in E$ with an ordered value domain,
    $\mathit{MinVal}_\rho : E \times T \rightarrow V$ returns the minimum value of $e$ up to and including time $t$:
    \[
        \mathit{MinVal}_\rho(e,\, t) \;=\; \inf\{\, \mathit{Val}_\rho(e,\, s) \mid s \in T,\; s \leq t \,\}
    \]
\end{definition}

The following two functions are defined for \textbf{discrete time only} ($T = T_D = \mathbb{N}$).

\begin{definition}[]\label{def:prev}
    $\mathit{Prev} : T_D \rightharpoonup T_D$ returns the preceding time step:
    \[
        \mathit{Prev}(t) \;=\; t - 1
    \]
    The function is partial: it is undefined at $t = 0$ since $t - 1 \notin T_D$.
\end{definition}

\begin{definition}[]\label{def:next}
    $\mathit{Next} : T_D \rightarrow T_D$ returns the following time step:
    \[
        \mathit{Next}(t) \;=\; t + 1
    \]
\end{definition}

\begin{definition}[]\label{def:reltime}
    $\mathit{RelTime} : T_D \times \mathbb{Z} \rightharpoonup T_D$ shifts a time point by $n$ steps, where $n \in \mathbb{Z}$ may be positive (forward) or negative (backward):
    \[
        \mathit{RelTime}(t,\, n) \;=\; t + n
    \]
    The function is partial: it is undefined when $t + n < 0$ since $t + n \notin T_D$.
    Note that $\mathit{Prev}(t) = \mathit{RelTime}(t, -1)$ and $\mathit{Next}(t) = \mathit{RelTime}(t, 1)$.
\end{definition}

\begin{definition}[]\label{def:diff}
    $\mathit{Diff} : T \times T \rightarrow \Delta$ returns the duration between two time points:
    \[
        \mathit{Diff}(t_1,\, t_2) \;=\; |t_1 - t_2|
    \]
\end{definition}

\begin{definition}[]\label{def:mk_interval}
    $\mathit{MkInterval} : T \times T \rightharpoonup \mathcal{I}$ constructs an interval from a start and end time point:
    \[
        \mathit{MkInterval}(t_1,\, t_2) \;=\; (t_1,\; t_2 - t_1)
    \]
    The function is partial: it is undefined when $t_1 > t_2$.
\end{definition}

\begin{definition}[Since-zero interval]\label{def:since_zero}
    $\mathit{SinceZero} : T \rightarrow \mathcal{I}$ constructs an interval from time zero to a given time point $t$:
    \[
        \mathit{SinceZero}(t) \;=\; (0,\; t)
    \]
\end{definition}

\begin{definition}[Interval start]\label{def:interval_start}
    $\mathit{Start} : \mathcal{I} \rightarrow T$ returns the start time point of an interval:
    \[
        \mathit{Start}((t,\, d)) \;=\; t
    \]
\end{definition}

\begin{definition}[Interval end]\label{def:interval_end}
    $\mathit{End} : \mathcal{I} \rightarrow T$ returns the end time point of an interval:
    \[
        \mathit{End}((t,\, d)) \;=\; t + d
    \]
\end{definition}

\begin{definition}[Interval duration]\label{def:interval_duration}
    $\mathit{Duration} : \mathcal{I} \rightarrow \Delta$ returns the duration of an interval:
    \[
        \mathit{Duration}((t,\, d)) \;=\; d
    \]
\end{definition}

\section{Operations}

Each entity type can be mapped to different \textit{operations}, which can be
\textit{executed} on an entity. Each \textit{operation} is mapped to a particular
predicate, which becomes \textit{true}, when the operation is \textit{executed}.

\begin{definition}[Operations set]\label{def:op_set}
    \[
        O = \{ op\,|\,\text{op is an operation causing an effect predicate to become true}\}
    \]
\end{definition}

\begin{definition}[Operations function]\label{def:op_fn}
    \[
        O_T: T_E \rightarrow O
    \]
\end{definition}

The effect predicate function $P_O$ maps each \textit{operation} to a predicate that becomes \text{true},
when \textit{executed}.

\begin{definition}[Effect predicate function]\label{def:ef_pred_fn}
    \[
        P_O: O \rightarrow \text{set of all possible predicates}
    \]
\end{definition}

\section{Notes}

\begin{itemize}
    \item Lepsi ponechat typ entity, protoze musime mapovat na modifiers,
          item ktere jsou rozdilne pro kazdy typ entity. Zaroven muzeme namapovat
          typ entity na mnozinu operaci.
    \item Modifiers need to be parameterized.
    \item We could theoretically map entity types directly to predicates, which
          can be set to \textit{true}, but maybe we will want different predicates depending on the context (e.g. discrete or continuous time).
    \item $T$ should be merged discrete and continuous.
    \item Need to resolve undefined with time functions.
    \item Last n-th occurence.
\end{itemize}

\end{document}